\documentclass[11pt]{article}

\usepackage[margin=1in]{geometry}

\title{Paired diagnostic report for RefCOCOg VLM grounding}
\author{}
\date{}

\begin{document}

\maketitle

\begin{abstract}
This report describes a paired diagnostic of four Qwen3-VL-8B-Instruct systems
for 2D RefCOCOg referring-expression grounding. The systems were evaluated on
the same 500 training examples and the same 200 validation examples using
corrected version 2 manifests. The validation set is a small diagnostic subset,
so these results are useful for this experiment but are not an official
benchmark or leaderboard result. SFT produced the largest validation gain over
the base model. Direct GRPO produced a smaller gain. GRPO after SFT was close
to SFT on validation and had the highest mean IoU on the 500-example training
set.
\end{abstract}

\section{Purpose and scope}

The repository studies referring-expression grounding with a vision-language
model. Given an image and a phrase describing one object, the model must return
one bounding box. The code uses Qwen3-VL-8B-Instruct and supports a base
evaluation, LoRA supervised fine-tuning (SFT), direct GRPO, and GRPO initialized
from an SFT adapter.

The experiment in this report answers a narrow question: how do these three
training variants compare with the base model when every system is evaluated
on the same examples and annotations? It does not estimate a leaderboard score,
generalization across RefCOCOg splits, or the effect of multiple random seeds.

\section{Data and paired evaluation}

The diagnostic used the following corrected manifests:

\begin{itemize}
  \item \texttt{refcocog\_train500\_v2.jsonl}, with 500 rows;
  \item \texttt{refcocog\_validation200\_v2.jsonl}, with 200 rows.
\end{itemize}

The two manifests have zero overlapping reference IDs. Within each split, all
four systems were evaluated on the same rows. The diagnostic script pairs rows
by \texttt{ref\_id} when available and checks the phrase and ground-truth pixel
box before calculating a comparison. The paired comparisons summarized here
were all comparable.

Earlier comparisons used mismatched manifests. Those comparisons were
superseded by this version 2 paired diagnostic and are not used in the results
below.

\section{Systems and training setup}

All systems use Qwen3-VL-8B-Instruct. The four configurations are:

\begin{description}
  \item[Base.] The base model without task training.
  \item[SFT.] LoRA supervised fine-tuning on the converted RefCOCOg examples.
  \item[Direct GRPO.] GRPO started from the base model.
  \item[SFT then GRPO.] GRPO started from the SFT adapter and used that adapter
  as the reference adapter.
\end{description}

The checked-in Slurm launcher uses bfloat16, LoRA rank 8, LoRA alpha 32,
\texttt{all-linear} target modules, maximum sequence length 2048, and one
example per device batch. SFT runs for one epoch with learning rate
\texttt{1e-4}, gradient accumulation 8, and the SFT JSONL data. The GRPO runs
use learning rate \texttt{5e-6}, four generations, \texttt{beta=0.04}, maximum
completion length 64, gradient accumulation 4, and the GRPO JSONL data.

The prompt asks for one box in the form
\texttt{[[x1, y1, x2, y2]]}, with coordinates on a 0 to 1000 scale relative
to the image width and height. The data converter stores the target box in
cached-image pixel coordinates for the reward function. The evaluator converts
valid normalized predictions to the same pixel space before computing IoU.

The GRPO reward is pixel-space IoU for a valid box and \texttt{-0.1} for an
invalid box. The evaluator records invalid or failed predictions as
\texttt{parse\_ok=false} and assigns them IoU 0.0.

\section{Metrics}

For each prediction file, the report uses:

\begin{description}
  \item[Mean IoU.] The mean intersection over union in cached-image pixel space.
  \item[Acc@0.5.] The fraction of examples with IoU at least 0.5.
  \item[Acc@0.75.] The fraction of examples with IoU at least 0.75.
  \item[Parse failures.] The number of examples without a valid parsed box.
\end{description}

The diagnostic also reports the mean predicted-area to ground-truth-area ratio
and mean center distance for parsed rows. The paired tables below focus on the
requested accuracy metrics and mean IoU. A paired mean-IoU delta is the variant
IoU minus the base IoU on each identical row, averaged over the split. The
improved, degraded, and unchanged counts classify those per-row deltas.

\section{Results}

\subsection{Validation results}

Table~\ref{tab:validation} reports the 200-example validation diagnostic.

\begin{table}[h]
\centering
\caption{Validation metrics on 200 paired examples.}
\label{tab:validation}
\begin{tabular}{lrrrr}
\hline
System & Mean IoU & Acc@0.5 & Acc@0.75 & Parse failures \\
\hline
Base & 0.8023357 & 0.845 & 0.780 & 1 \\
SFT & 0.8160114 & 0.860 & 0.800 & 0 \\
Direct GRPO & 0.8042084 & 0.850 & 0.780 & 1 \\
SFT then GRPO & 0.8158485 & 0.860 & 0.795 & 0 \\
\hline
\end{tabular}
\end{table}

Relative to the base validation run, the paired mean-IoU deltas were:

\begin{table}[h]
\centering
\caption{Paired validation mean-IoU changes against the base run.}
\label{tab:validation-paired}
\begin{tabular}{lrrrrr}
\hline
Variant & Delta & Improved & Degraded & Unchanged & Parse failures \\
\hline
SFT & +0.01367575 & 101 & 97 & 2 & 0 \\
Direct GRPO & +0.00187275 & 44 & 44 & 112 & 1 \\
SFT then GRPO & +0.01351283 & 99 & 99 & 2 & 0 \\
\hline
\end{tabular}
\end{table}

SFT has the highest validation mean IoU and Acc@0.75 among the trained
variants. SFT then GRPO has the same Acc@0.5 as SFT and a slightly lower mean
IoU and Acc@0.75 on this split. Direct GRPO is close to the base model: its
mean IoU is higher by 0.00187275, its Acc@0.5 is higher by 0.005, and its
Acc@0.75 is unchanged. The base and direct GRPO runs each have one parse
failure. SFT and SFT then GRPO have none.

\subsection{Training-split results}

Table~\ref{tab:training} reports the same systems on the 500-example training
diagnostic.

\begin{table}[h]
\centering
\caption{Training-split metrics on 500 paired examples.}
\label{tab:training}
\begin{tabular}{lrrrr}
\hline
System & Mean IoU & Acc@0.5 & Acc@0.75 & Parse failures \\
\hline
Base & 0.7976700 & 0.858 & 0.744 & 0 \\
SFT & 0.8251656 & 0.890 & 0.782 & 0 \\
Direct GRPO & 0.8019457 & 0.864 & 0.748 & 0 \\
SFT then GRPO & 0.8291169 & 0.890 & 0.794 & 0 \\
\hline
\end{tabular}
\end{table}

The paired training-split changes were:

\begin{table}[h]
\centering
\caption{Paired training-split mean-IoU changes against the base run.}
\label{tab:training-paired}
\begin{tabular}{lrrrr}
\hline
Variant & Delta & Improved & Degraded & Unchanged \\
\hline
SFT & +0.02749562 & 258 & 237 & 5 \\
Direct GRPO & +0.00427569 & 105 & 96 & 299 \\
SFT then GRPO & +0.03144688 & 278 & 216 & 6 \\
\hline
\end{tabular}
\end{table}

SFT then GRPO has the highest training mean IoU and Acc@0.75. SFT and SFT
then GRPO both reach Acc@0.5 of 0.890. Direct GRPO again has a smaller change
from base than SFT, with a paired mean-IoU delta of +0.00427569.

\subsection{Training run records}

The recorded training runs completed the following steps:

\begin{table}[h]
\centering
\caption{Training run records. These values are not comparable across methods.}
\label{tab:training-runs}
\begin{tabular}{lrrll}
\hline
Run & Steps & Runtime (s) & Train loss & Final logged value \\
\hline
SFT & 63/63 & 130.4228 & 0.89258396 & token accuracy 0.69059829 \\
Direct GRPO & 500/500 & 2134.9276 & 0.00159036 & reward 0.68845888 \\
SFT then GRPO & 500/500 & 2531.4886 & 0.00034432 & reward 0.69928728 \\
\hline
\end{tabular}
\end{table}

The SFT record contains token accuracy, while the GRPO records contain a final
logged reward. The objectives, logging values, and training procedures differ,
so the train losses and final values should not be ranked across methods.

\section{Analysis}

The paired results support three limited observations.

First, SFT gives the clearest improvement over the base model in both measured
splits. Its mean-IoU delta is +0.01367575 on validation and +0.02749562 on the
training diagnostic. The validation Acc@0.5 rises from 0.845 to 0.860, and
Acc@0.75 rises from 0.780 to 0.800.

Second, direct GRPO gives a smaller change in this run. Its validation paired
delta is +0.00187275, with 44 improved rows, 44 degraded rows, and 112
unchanged rows. The training paired delta is +0.00427569, with 105 improved,
96 degraded, and 299 unchanged rows. These results do not show a large direct
GRPO improvement over the base model on the evaluated examples.

Third, SFT then GRPO remains close to SFT on validation and exceeds it on the
training diagnostic. On validation, SFT then GRPO has mean IoU 0.8158485,
compared with 0.8160114 for SFT. On the training diagnostic, it has mean IoU
0.8291169, compared with 0.8251656 for SFT. This difference is descriptive;
the experiment does not include repeated runs or a larger held-out benchmark
that would establish whether the difference is reliable.

The parse-failure counts are low. Validation has one failure for the base and
direct GRPO runs and none for SFT or SFT then GRPO. All four training-split
runs have zero parse failures. Because invalid parses receive IoU 0, parse
reliability is part of the reported mean, not a separate post-processing
adjustment.

\section{Limitations}

The validation diagnostic contains only 200 examples. It is not a definitive
benchmark, and the results should not be presented as an official RefCOCOg
leaderboard result. The training split is used here as a diagnostic comparison,
not as evidence of held-out generalization.

The experiment reports one run for each configuration. It does not measure
variation across random seeds, compare multiple hyperparameter settings, or
provide statistical uncertainty. The paired comparison reduces variation from
example composition, but it does not remove these other sources of uncertainty.

The report includes mean IoU, threshold accuracies, parse failures, and paired
win/loss counts. The diagnostic JSON also contains area-ratio and center-distance
statistics. These box-shape measures are useful for follow-up analysis, but
they do not replace a larger evaluation.

\section{Artifacts and reproduction}

The completed diagnostic artifacts remain on the HPC filesystem and are not
tracked in this repository:

\begin{verbatim}
/scratch/ambpdc/soham/vlm/outputs/diagnostic/
base_train.jsonl
sft_train.jsonl
grpo_train_h100.jsonl
sft_grpo_train_h100.jsonl
base_val.jsonl
sft_val.jsonl
grpo_val.jsonl
sft_grpo_val.jsonl
full_report.json
\end{verbatim}

The report can be reproduced from the repository with the following workflow.
The commands assume the HPC layout used by the Slurm launcher.

\begin{enumerate}
  \item Install the lightweight dependencies from \texttt{requirements.txt}.
  \item Prepare the RefCOCOg image caches and the version 2 train and validation
  manifests with \texttt{scripts/prepare\_refcocog\_images.py}.
  \item Convert each manifest with
  \texttt{scripts/prepare\_swift\_data.py}.
  \item Run the three training modes from the repository checkout:
  \texttt{sft}, \texttt{grpo\_base}, and \texttt{grpo\_after\_sft}.
  \item Evaluate the base model and each adapter with
  \texttt{scripts/evaluate\_grounding.py}, using the same manifest within each
  split.
  \item Run \texttt{scripts/analyze\_grounding\_diagnostic.py} with the eight
  prediction files and write \texttt{full\_report.json}.
\end{enumerate}

The paired diagnostic command is:

\begin{verbatim}
python3 scripts/analyze_grounding_diagnostic.py \
  --prediction base/train=/scratch/ambpdc/soham/vlm/outputs/diagnostic/base_train.jsonl \
  --prediction sft/train=/scratch/ambpdc/soham/vlm/outputs/diagnostic/sft_train.jsonl \
  --prediction grpo/train=/scratch/ambpdc/soham/vlm/outputs/diagnostic/grpo_train_h100.jsonl \
  --prediction sft_grpo/train=/scratch/ambpdc/soham/vlm/outputs/diagnostic/sft_grpo_train_h100.jsonl \
  --prediction base/val=/scratch/ambpdc/soham/vlm/outputs/diagnostic/base_val.jsonl \
  --prediction sft/val=/scratch/ambpdc/soham/vlm/outputs/diagnostic/sft_val.jsonl \
  --prediction grpo/val=/scratch/ambpdc/soham/vlm/outputs/diagnostic/grpo_val.jsonl \
  --prediction sft_grpo/val=/scratch/ambpdc/soham/vlm/outputs/diagnostic/sft_grpo_val.jsonl \
  --output /scratch/ambpdc/soham/vlm/outputs/diagnostic/full_report.json
\end{verbatim}

Model weights, images, manifests, adapters, logs, and prediction files are
external artifacts. Their absence from Git is intentional.

\section{Conclusion}

On this paired diagnostic, SFT produced the largest validation improvement over
the base model. Direct GRPO produced a smaller improvement. SFT then GRPO was
very close to SFT on validation and had the strongest training-split score.
These findings describe the measured 500 and 200 example diagnostics only.

\end{document}
